\documentclass[10pt]{article}
\usepackage[T1]{fontenc}
\usepackage[utf8]{inputenc}
\usepackage[french]{babel}
\usepackage{amsmath,amssymb,mathrsfs,mathpazo}
\usepackage{geometry}
\geometry{a4paper,margin=2cm}
\usepackage{graphicx,array,tabularx,tabularray,colortbl,multirow,multicol,booktabs,enumitem,pifont,ulem,fancybox,diagbox,cancel}
\usepackage{pstricks,pstricks-add,pst-plot,pst-tree,pst-node,pst-text,pst-func,pst-math,pst-bspline,pst-3dplot,pst-eucl,pst-solides3d}
\let\pscalculate\undefined
\usepackage{pst-bezier}
\usepackage{tikz,pgfplots,tkz-tab}
\usetikzlibrary{arrows,arrows.meta,patterns,calc,positioning,trees,decorations.markings,shapes.geometric,shapes.misc,shapes.arrows,intersections,decorations.pathmorphing,decorations.pathreplacing,backgrounds,shadows,babel}
\pgfplotsset{compat=1.18}
\usepgfplotslibrary{fillbetween}
\usepackage[np]{numprint}
\usepackage{listings,makecell,pas-tableur,twoopt}
\let\barre\undefined
\usepackage[pstricks]{tablvar}
\definecolor{darkgreen}{RGB}{0,100,0}
\providecommand{\red}{\color{red}}
\providecommand{\blue}{\color{blue}}
\providecommand{\green}{\color{green}}
\providecommand{\black}{\color{black}}
\providecommand{\white}{\color{white}}
\providecommand{\gray}{\color{gray}}
\providecommand{\cyan}{\color{cyan}}
\providecommand{\magenta}{\color{magenta}}
\providecommand{\yellow}{\color{yellow}}
\providecommand{\violet}{\color{violet}}
\providecommand{\orange}{\color{orange}}
\providecommand{\euro}{\texteuro}
\providecommand{\EUR}{\texteuro}
\providecommand{\ieme}{\textsuperscript{e}}
\providecommand{\degres}{\ensuremath{^\circ}}
\providecommand{\degree}{\ensuremath{^\circ}}
\providecommand{\vv}[1]{\overrightarrow{#1}}
\providecommand{\Coord}[1]{#1}
\providecommand{\decosix}{}
\providecommand{\starredbullet}{\ensuremath{\bullet}}
\pagestyle{empty}
\setlength{\parindent}{0pt}
\usepackage{adjustbox}\DeclareUnicodeCharacter{2236}{:}\DeclareUnicodeCharacter{2190}{\ensuremath{\leftarrow}}
\providecommand{\R}{\mathbb{R}}
\providecommand{\N}{\mathbb{N}}
\providecommand{\Z}{\mathbb{Z}}
\providecommand{\Q}{\mathbb{Q}}
\providecommand{\D}{\mathbb{D}}
\providecommand{\C}{\mathbb{C}}
\providecommand{\e}{{\mathrm{e}}}
\providecommand{\vect}[1]{\overrightarrow{#1}}
\providecommand{\vectt}[1]{\overrightarrow{\mathrm{#1}}}
\providecommand{\vvt}[1]{\overrightarrow{\mathrm{#1}}}
\providecommand{\np}[2][]{{#2~#1}}
\providecommand{\vv}[1]{{\overrightarrow{#1}}}
\renewcommand{\d}{{\mathrm{d}}}
\providecommand{\degres}{{^\circ}}
\providecommand{\degre}{{^\circ}}
\providecommand{\blue}{\color{blue}}
\providecommand{\red}{\color{red}}
\providecommand{\enskip}{\quad}
\providecommand{\llbracket}{\lbrack\!\lbrack}
\providecommand{\rrbracket}{\rbrack\!\rbrack}
\providecommand{\euro}{{\text{€}}}
\providecommand{\no}{{\text{n°}}}
\providecommand{\interval}[2]{{[#1;#2]}}
\providecommand{\intervFO}[2]{{[#1;#2[}}
\providecommand{\dotfill}{\ldots}

\providecommand{\ds}{\displaystyle}
\providecommand{\textding}[1]{#1}
\providecommand{\Oij}{\(\mathrm{O};\vec{\imath},\vec{\jmath}\)}
\providecommand{\Oijk}{\(\mathrm{O};\vec{\imath},\vec{\jmath},\vec{k}\)}
\providecommand{\Ouv}{\(\mathrm{O};\vec{u},\vec{v}\)}
\providecommand{\intervFF}[2]{\ensuremath{[#1;#2]}}
\renewcommand{\euro}{\ifmmode\text{\texteuro}\else\texteuro\fi}
\DeclareUnicodeCharacter{00B0}{\ifmmode^{\circ}\else\textdegree\fi}
\DeclareUnicodeCharacter{20AC}{\ifmmode\text{\texteuro}\else\texteuro\fi}
% Énoncés : Éducation nationale. Transcription et corrigés : APMEP, Denis Vergès et contributeurs. Découpage et classement : CoopMaths, Sébastien Lozano, Éric Elter et contributeurs. Adaptation de la présentation : Maths Explorer.
% CC BY-SA : https://coopmaths.fr/www/about/
% https://coopmaths.fr/alea/static/bac/2021/tex/bac_2021_05_ameriquenord_1.tex
\begin{document}
\noindent\textit{APMEP / CoopMaths. CC BY-SA. Présentation adaptée par Maths Explorer.}\par\medskip
\textbf{Commun à tous les candidats}

\medskip


Les probabilités demandées dans cet exercice seront arrondies à $10^{-3}$.

\medskip

Un laboratoire pharmaceutique vient d'élaborer un nouveau test anti-dopage. 

\bigskip

\textbf{Partie A}

\medskip

Une étude sur ce nouveau test donne les résultats suivants:

\setlength\parindent{9mm}
\begin{itemize}
\item[$\bullet~~$] si un athlète est dopé, la probabilité que le résultat du test soit positif est $0,98$ (sensibilité du test) ;
\item[$\bullet~~$]si un athlète n'est pas dopé, la probabilité que le résultat du test soit négatif est $0,995$ (spécificité du test).
\end{itemize}
\setlength\parindent{0mm}

\smallskip

On fait subir le test à un athlète sélectionné au hasard au sein des participants à une compétition d'athlétisme. 

On note $D$ l'évènement \og l'athlète est dopé \fg{} et $T$ l'évènement \og le test est positif \fg. 

On admet que la probabilité de l'évènement $D$ est égale à 0,08.

\medskip

\begin{enumerate}
\item Traduire la situation sous la forme d'un arbre pondéré.
\item Démontrer que $P(T) = 0,083$.
\item 
	\begin{enumerate}
		\item Sachant qu'un athlète présente un test positif, quelle est la probabilité qu'il soit dopé ?
		\item Le laboratoire décide de commercialiser le test si la probabilité de l'évènement \og un athlète présentant un test positif est dopé \fg{} est supérieure ou égale à $0,95$.

Le test proposé par le laboratoire sera-t-il commercialisé ? Justifier.
	\end{enumerate}
\end{enumerate}

\bigskip

\textbf{Partie B}

\medskip

Dans une compétition sportive, on admet que la probabilité qu'un athlète contrôlé présente un test positif est $0,103$.

\medskip

\begin{enumerate}
\item Dans cette question \textbf{1.} on suppose que les organisateurs décident de contrôler $5$ ~athlètes au hasard parmi les athlètes de cette compétition. 

On note $X$ la variable aléatoire égale au nombre d'athlètes présentant un test positif parmi les $5$ athlètes contrôlés.
	\begin{enumerate}
		\item Donner la loi suivie par la variable aléatoire $X$. Préciser ses paramètres.
		\item Calculer l'espérance $E(X)$ et interpréter le résultat dans le contexte de l'exercice.
		\item Quelle est la probabilité qu'au moins un des 5 athlètes contrôlés présente un test positif ?
	\end{enumerate}
\item Combien d'athlètes faut-il contrôler au minimum pour que la probabilité de l'évènement \og au moins un athlète contrôlé présente un test positif\fg{} soit supérieure ou égale à $0,75$ ? Justifier.
\end{enumerate}
\vspace{0,5cm}
\bigskip

\end{document}
